\section{Order forcing}

\begin{theorem}[Binary-receipt order forcing]
For a deterministic \(C_2\)-covariant two-sheeted lift of \(C_4\):
\begin{enumerate}
\item \(\omega=0\) gives \(C_4\sqcup C_4\) and successor order
four;
\item \(\omega=1\) gives \(C_8\) and successor order eight;
\item two voltage assignments are gauge equivalent exactly when they
have the same holonomy.
\end{enumerate}
\end{theorem}

\begin{proof}
The four-step law is \(T^4(i,s)=(i,s+\omega)\).
For \(\omega=0\), each sheet closes separately after four steps.
For \(\omega=1\), \(T^4\) swaps sheets, so \(T^4\neq\id\) but
\(T^8=\id\). Projection has order four, hence \(T\) has order
eight and one orbit on all eight states.

A gauge \(g:V\to\Ztwo\) changes voltages by
\[
v_i'=v_i+g_i+g_{i+1},
\]
which preserves their sum. Conversely, if \(v\) and \(v'\) have
equal sums, set \(d_i=v_i+v_i'\), choose \(g_0=0\), and define
\(g_{i+1}=g_i+d_i\). The zero-sum condition closes this recursion,
so equal holonomy implies gauge equivalence.
\end{proof}

\begin{corollary}
A visible four-cycle with binary fiber and nontrivial circuit receipt
has a unique connected lift up to fiber-preserving isomorphism:
\(C_8\).
\end{corollary}

In the nontrivial class,
\[
pT^4=p,\qquad T^4\neq\id,\qquad T^8=\id.
\]
Visible return therefore precedes lifted closure.
